% ==========================================
% BAB V IMPLEMENTASI
% ==========================================
\cleardoublepage
\chapter{IMPLEMENTASI}
\label{chap:implementasi}

\section{Gambaran Umum Implementasi}
\label{sec:gambaran-umum-implementasi}

Bab ini menjelaskan realisasi rancangan arsitektur sistem inspeksi kontainer digital
terintegrasi yang telah dibahas pada Bab IV. Fokus pembahasan tidak diarahkan
pada rincian implementasi kode program secara umum, melainkan pada bagaimana
keputusan arsitektural direalisasikan menjadi sistem yang dapat beroperasi.
Dengan demikian, bab ini menempatkan implementasi sebagai bukti bahwa
rancangan arsitektur yang diusulkan dapat diwujudkan menjadi solusi yang
terintegrasi.

Dengan posisi tersebut, pembahasan pada bab ini tidak dimaksudkan sebagai
dokumentasi menyeluruh atas struktur kode sumber atau keputusan pemrograman
setiap komponen. Implementasi dibaca sebagai konteks realisasi yang
menunjukkan apakah rancangan komponen, kontrak data, dan alur integrasi yang
diusulkan pada tugas akhir ini tercermin pada sistem yang dibangun.

Artefak sistem yang digunakan untuk merealisasikan rancangan terdiri atas beberapa komponen utama yang
dipisahkan berdasarkan tanggung jawabnya. Pemisahan tersebut selaras dengan
prinsip modularitas, pemrosesan tersebar, dan integrasi data yang telah
ditetapkan pada tahap perancangan. Secara umum, sistem terdiri atas komponen
\textit{edge} untuk pengolahan aliran video dan inferensi awal, komponen
layanan aplikasi untuk orkestrasi proses inspeksi serta persistensi data, komponen
aplikasi web untuk pemantauan dan inspeksi, serta komponen penyimpanan data
untuk informasi terstruktur dan artefak visual hasil inspeksi.

Realisasi ini menghasilkan dua jalur utama pertukaran data. Jalur pertama adalah
jalur data terstruktur yang mengalir dari komponen \textit{edge} ke layanan API,
kemudian disimpan pada basis data dan diakses oleh aplikasi web. Jalur kedua
adalah jalur video langsung yang mengalir dari komponen \textit{edge} ke aplikasi
web melalui \textit{WebSocket}. Pemisahan dua jalur ini merupakan keputusan
arsitektural penting karena data transaksional dan aliran video memiliki
karakteristik kebutuhan yang berbeda dari sisi latensi, beban jaringan, dan pola
akses.

Karena kontribusi tugas akhir ini berada pada perancangan arsitektur sistem,
ruang lingkup realisasi perlu dibedakan dari rancangan sasaran yang lebih luas.
Tabel~\ref{tbl:batas-realisasi-sasaran} menunjukkan bagian yang telah
direalisasikan sebagai bukti implementabilitas rancangan dan bagian yang tetap
diposisikan sebagai kesiapan pengembangan lanjutan.

\begin{table}[ht]
  \centering
  \footnotesize
  \caption[Batas realisasi dan rancangan sasaran]{Batas Realisasi dan Rancangan Sasaran}
  \label{tbl:batas-realisasi-sasaran}
  \setlength{\tabcolsep}{4pt}
  \renewcommand{\arraystretch}{1.12}
  \begin{tabularx}{\textwidth}{>{\raggedright\arraybackslash}p{0.24\textwidth} >{\raggedright\arraybackslash}X >{\raggedright\arraybackslash}X}
    \toprule
    \textbf{Aspek} & \textbf{Direalisasikan dalam Tugas Akhir} & \textbf{Diposisikan sebagai Rancangan Sasaran} \\
    \midrule
    Integrasi internal &
    Komunikasi komponen \textit{edge}, layanan API, aplikasi web, MongoDB, dan Cloudflare R2 pada alur inspeksi utama &
    Integrasi operasional penuh dengan TOS, PCS, INSW, atau sistem eksternal pelabuhan lain. \\
    Model data inspeksi &
    Agregasi hasil OCR, hasil pemindaian, konteks manifes, status, dan lampiran bukti pada entitas kontainer &
    Pengayaan data dari seluruh sumber eksternal pelabuhan dan standardisasi lintas organisasi. \\
    Jalur video langsung &
    Penyiaran \textit{frame} beranotasi dari komponen \textit{edge} ke aplikasi web melalui \textit{WebSocket} &
    Infrastruktur produksi multi-lokasi dengan jaminan ketersediaan dan pengukuran performa skala operasional. \\
    Artefak bukti inspeksi &
    Penyimpanan citra hasil OCR, visualisasi deteksi, dan bukti pendukung pada penyimpanan objek &
    Kebijakan retensi, tata kelola audit formal, dan integrasi bukti ke sistem kepabeanan atau terminal. \\
    Evaluasi &
    Validasi skenario arsitektur: komunikasi antarkomponen, konsistensi data, pemisahan jalur, dan ketertelusuran bukti &
    Evaluasi bisnis, evaluasi \textit{deployment} produksi, dan evaluasi kuantitatif mendalam terhadap performa model AI. \\
    \bottomrule
  \end{tabularx}
\end{table}

\FloatBarrier

\section{Realisasi Arsitektur Sistem}
\label{sec:realisasi-arsitektur-sistem}

\subsection{Struktur Komponen Implementasi}
\label{subsec:struktur-komponen-implementasi}

Realisasi sistem memperlihatkan empat kelompok komponen utama yang saling
terhubung.

\begin{enumerate}
  \item Komponen \textit{edge} direalisasikan sebagai layanan berbasis Python dan
  \textit{FastAPI}. Komponen ini bertugas menarik aliran video dari kamera
  \textit{RTSP}, menjalankan model deteksi dan \textit{optical character
  recognition} (OCR), mengelola aliran video, serta mengirimkan hasil pengolahan
  ke komponen lain. Pada tahap implementasi, komponen \textit{edge} juga
  menyediakan layanan \textit{WebSocket} untuk menyiarkan \textit{frame} yang
  sudah dianotasi kepada aplikasi web.

  \item Komponen layanan API direalisasikan sebagai layanan aplikasi berbasis
  Node.js dan Express. Komponen ini berperan sebagai pusat orkestrasi data
  terstruktur, autentikasi, pengelolaan status inspeksi, serta antarmuka layanan
  bagi aplikasi web dan komponen \textit{edge}. Pada operasi inti, seluruh hasil
  identifikasi kontainer dan hasil pemindaian dari komponen \textit{edge}
  diterima oleh layanan ini melalui keluarga \textit{endpoint} \texttt{/api/edge}.
  Kanal pelacakan tambahan yang sempat tersedia tidak diposisikan sebagai jalur
  utama pada arsitektur inspeksi yang dibahas dalam tugas akhir ini.

  \item Komponen aplikasi web direalisasikan sebagai aplikasi dasbor berbasis
  web. Komponen ini menyediakan antarmuka bagi operator untuk memantau aliran
  video, meninjau hasil inspeksi, mengelola aliran kamera, dan membaca hasil
  identifikasi serta analisis yang sudah tersimpan. Aplikasi web mengakses data
  terstruktur melalui layanan API dan menerima video langsung dari komponen
  \textit{edge} melalui \textit{WebSocket}. Penyebutan teknologi implementasi web
  tidak menjadi fokus utama pada tugas akhir ini, karena yang ditekankan adalah
  perannya sebagai lapisan presentasi dalam arsitektur sistem.

  \item Komponen penyimpanan direalisasikan dalam dua bentuk, yaitu MongoDB
  sebagai basis data utama untuk data terstruktur dan Cloudflare R2 sebagai
  media penyimpanan artefak visual seperti citra hasil identifikasi, citra hasil
  pemindaian, dan visualisasi anotasi. Pemisahan ini sesuai dengan prinsip bahwa
  data terstruktur dan data media memiliki pola penyimpanan yang berbeda.
\end{enumerate}

\subsection{Pemetaan Implementasi terhadap Lapisan Arsitektur}
\label{subsec:pemetaan-implementasi-lapisan}

Hasil realisasi sistem menunjukkan bahwa struktur sistem yang dibangun selaras
dengan arsitektur berlapis yang telah dirancang sebelumnya. \textit{Presentation
Layer} direalisasikan oleh aplikasi web yang menyediakan fitur pemantauan,
inspeksi, pelaporan, dan pengelolaan aliran kamera. \textit{Application Layer}
direalisasikan oleh layanan API yang mengelola logika proses inspeksi,
validasi data, autentikasi, dan orkestrasi layanan. \textit{Data Layer}
direalisasikan oleh MongoDB sebagai penyimpanan data terstruktur dan
Cloudflare R2 sebagai penyimpanan artefak inspeksi berbasis berkas. Adapun
\textit{Edge Layer} direalisasikan oleh layanan \textit{edge} yang
berinteraksi langsung dengan kamera, menjalankan inferensi, dan mengirimkan
hasil pengolahan ke komponen lain.

Pemetaan tersebut menunjukkan bahwa realisasi sistem tidak
meninggalkan struktur rancangan konseptual, melainkan menerjemahkannya menjadi
komponen nyata yang dapat diuji. Dengan demikian, hasil realisasi mendukung
argumen bahwa rancangan arsitektur yang disusun pada tugas akhir ini bersifat
dapat diimplementasikan.

Untuk memperjelas hubungan antara rancangan pada Bab IV dan realisasi sistem,
Tabel~\ref{tbl:pemetaan-rancangan-realisasi} merangkum komponen rancangan,
bentuk realisasi, bukti atau penanda implementasi, serta implikasinya terhadap
arsitektur.

\begin{table}[htbp]
  \centering
  \caption{Pemetaan Rancangan Arsitektur terhadap Realisasi Sistem}
  \label{tbl:pemetaan-rancangan-realisasi}
  \small
  \setlength{\tabcolsep}{4pt}
  \renewcommand{\arraystretch}{1.08}
  \begin{tabularx}{\textwidth}{>{\raggedright\arraybackslash}p{0.18\textwidth} >{\raggedright\arraybackslash}p{0.24\textwidth} >{\raggedright\arraybackslash}p{0.24\textwidth} >{\raggedright\arraybackslash}X}
    \toprule
    \textbf{Rancangan} & \textbf{Realisasi} & \textbf{Bukti/Penanda} & \textbf{Implikasi Arsitektural} \\
    \midrule
    \textit{Presentation Layer} &
    Aplikasi web untuk pemantauan, peninjauan hasil, dan pelaporan inspeksi &
    Halaman pemantauan dan halaman detail kontainer &
    Hasil inspeksi dapat dibaca operator melalui antarmuka yang konsisten. \\
    \textit{Application Layer} &
    Layanan API sebagai titik orkestrasi dan persistensi data &
    Kontrak HTTP pada keluarga \textit{endpoint} \texttt{/api/edge/*} &
    Data inspeksi masuk melalui jalur terstruktur dan dapat divalidasi sebelum disimpan. \\
    \textit{Data Layer} &
    MongoDB untuk data terstruktur dan Cloudflare R2 untuk artefak visual &
    Dokumen kontainer, URL citra, dan bukti visual hasil pemindaian &
    Data transaksional dan media inspeksi dipisahkan sesuai karakteristik penyimpanannya. \\
    \textit{Edge Layer} &
    Layanan \textit{edge} untuk akuisisi video, inferensi awal, dan siaran langsung &
    Aliran \textit{RTSP}, hasil OCR, deteksi, dan \textit{WebSocket} &
    Pemrosesan yang sensitif terhadap latensi dilakukan dekat sumber data. \\
    Propagasi konteks inspeksi &
    Nomor kontainer hasil OCR digunakan oleh alur pemindaian lain &
    Hasil OCR, hasil kerusakan, barang berisiko, dan kategori berada pada entitas kontainer yang sama &
    Integrasi lintas alur dapat dilakukan tanpa pengisian ulang identitas kontainer. \\
    \bottomrule
  \end{tabularx}
\end{table}

\FloatBarrier

\section{Realisasi Alur Data dan Integrasi}
\label{sec:realisasi-alur-data}

\subsection{Jalur Data Terstruktur}
\label{subsec:jalur-data-terstruktur}

Jalur data terstruktur merupakan jalur utama untuk menyimpan hasil inspeksi dan
menyediakan data operasional kepada pengguna. Pada implementasinya, komponen
\textit{edge} mengirimkan hasil identifikasi kontainer, hasil pemindaian
kerusakan, hasil pemindaian barang berisiko, dan hasil klasifikasi muatan ke
layanan API melalui antarmuka HTTP pada keluarga \textit{endpoint} \texttt{/api/edge/*}.
Dengan demikian, jalur operasi utama sistem terpusat pada satu kontrak layanan
yang konsisten untuk hasil inspeksi kontainer.

Alur tersebut dimulai ketika kamera mengirimkan aliran video ke komponen
\textit{edge}. Komponen \textit{edge} kemudian memproses \textit{frame},
mengubah hasil inferensi menjadi data terstruktur, dan mengirimkannya ke
layanan API. Selanjutnya, layanan API memvalidasi, menyimpan, dan memperbarui
data pada basis data, sedangkan aplikasi web membaca data tersebut melalui
antarmuka HTTP yang sama.

Dengan pola tersebut, aplikasi web tidak berinteraksi langsung dengan logika
inferensi pada komponen \textit{edge} untuk memperoleh hasil inspeksi. Seluruh
data operasional yang bersifat persisten tetap melalui layanan API sehingga
konsistensi data dan kontrol akses dapat dijaga secara terpusat. Jalur
\texttt{/api/tracking-data} yang masih tersedia pada implementasi tidak menjadi
jalur inti untuk operasi inspeksi kontainer, sehingga tidak dijadikan dasar
utama pembahasan pada naskah ini.

\subsection{Jalur Video Langsung}
\label{subsec:jalur-video-langsung}

Selain jalur data terstruktur, sistem juga merealisasikan jalur video langsung
antara komponen \textit{edge} dan aplikasi web. Jalur ini digunakan untuk
menyajikan \textit{frame} yang telah dianotasi secara waktu nyata. Pemisahan
jalur ini dilakukan agar layanan API tidak dibebani oleh arus video berukuran
besar yang tidak memerlukan penyimpanan transaksional.

Pada implementasinya, aplikasi web membuka koneksi \textit{WebSocket} langsung ke
komponen \textit{edge}. Melalui koneksi tersebut, aplikasi web menerima
\textit{frame} yang sudah diberi anotasi beserta \textit{metadata} deteksi yang relevan
untuk pemantauan. Dengan pendekatan ini, operator dapat memantau kondisi
lapangan secara langsung, tetapi penyimpanan hasil inspeksi formal tetap
dilakukan melalui layanan API. Keputusan ini memperlihatkan pemisahan yang
jelas antara kebutuhan pemantauan waktu nyata dan kebutuhan persistensi
data.

Hubungan antara jalur data terstruktur dan jalur video langsung dirangkum pada
Gambar~\ref{fig:implementation_data_flow}. Diagram tersebut memperlihatkan
bahwa layanan API menjadi pusat pengelolaan data persisten, sedangkan aliran
video langsung tetap bergerak dari komponen \textit{edge} ke aplikasi web tanpa
melalui jalur penyimpanan transaksional.

\begin{figure}[htbp]
  \centering
  \resizebox{\textwidth}{!}{%
  \begin{tikzpicture}[
    scale=0.92,
    transform shape,
    block/.style={draw, rounded corners, align=center, minimum width=2.45cm, minimum height=0.9cm, font=\small},
    data/.style={draw, align=center, minimum width=2.45cm, minimum height=0.9cm, font=\small},
    arrow/.style={-{Latex[length=2mm]}, thick}
  ]
    \node[block] (camera) at (0,2.2) {Kamera\\\textit{RTSP}};
    \node[block] (tepi) at (3.4,2.2) {Komponen\\\textit{Edge}};
    \node[block] (api) at (6.8,2.2) {Layanan API\\\texttt{/api/...}};
    \node[data] (db) at (10.2,2.2) {MongoDB\\\textit{containers}};
    \node[block] (web) at (3.4,-0.35) {Aplikasi Web};
    \node[data] (r2) at (8.2,-0.35) {Cloudflare R2\\Artefak Visual};

    \draw[arrow] (camera) -- (tepi);
    \draw[arrow] (tepi) -- (api);
    \draw[arrow] (api) -- (db);
    \draw[arrow] (api.south east) -- node[right,font=\scriptsize]{simpan artefak} (r2.north);
    \draw[arrow] (web.north east) -- node[left,font=\scriptsize]{REST} (api.south west);
    \draw[arrow] (web) -- node[below,font=\scriptsize]{akses URL} (r2);
    \draw[arrow] (tepi.south) -- node[left,font=\scriptsize]{\textit{WebSocket}} (web.north);
  \end{tikzpicture}}
  \caption{Alur Realisasi Data Inspeksi pada Implementasi Sistem}
  \label{fig:implementation_data_flow}
\end{figure}

\FloatBarrier

\subsection{Propagasi Nomor Kontainer}
\label{subsec:propagasi-nomor-kontainer}

Salah satu kebutuhan integrasi utama dalam sistem ini adalah kemampuan
menyelaraskan hasil identifikasi kontainer dengan aliran pemeriksaan lainnya.
Pada realisasinya, ketika aliran OCR berhasil mendeteksi nomor kontainer,
nomor tersebut secara otomatis dipropagasikan ke aliran lain yang sedang aktif.

Mekanisme ini memiliki dua implikasi arsitektural yang penting:

\begin{enumerate}
  \item Hasil OCR tidak berhenti sebagai keluaran lokal pada satu model,
  melainkan menjadi konteks bersama bagi proses inspeksi lanjutan.
  \item Hasil pemindaian kerusakan, barang berisiko, dan kategori muatan dapat
  dikaitkan ke entitas kontainer yang sama tanpa entri manual berulang dari
  operator.
\end{enumerate}

Dengan demikian, realisasi sistem berhasil menunjukkan kebutuhan integrasi data
lintas modul yang sebelumnya didefinisikan sebagai salah satu kapabilitas
arsitektur inti.

\section{Realisasi Arsitektur Data}
\label{sec:realisasi-arsitektur-data}

\subsection{Struktur Data Inti}
\label{subsec:struktur-data-inti}

Fokus utama kontribusi pada tugas akhir ini berada pada bagaimana data inspeksi
dimodelkan, disimpan, dan dipertukarkan secara konsisten. Realisasi sistem
menunjukkan bahwa operasi inspeksi utama direalisasikan dengan menempatkan
\textit{Container} sebagai pusat agregasi hasil inspeksi, \textit{Manifest}
sebagai konteks deklarasi muatan, dan \textit{User} sebagai konteks autentikasi
serta akses pengguna.

\textbf{Entitas \textit{Container}}

Entitas \textit{Container} menjadi pusat agregasi data inspeksi. Entitas ini
menyimpan identitas kontainer, hasil OCR, ringkasan manifest, hasil pemindaian
kerusakan, hasil pemindaian barang berisiko, hasil klasifikasi muatan, keputusan
sistem, hasil tinjauan manual, serta lampiran bukti inspeksi. Secara
implementatif, struktur ini direalisasikan melalui objek-objek tertanam seperti
\textit{ocr}, \textit{manifest}, \textit{defectScan}, \textit{illegalScan},
\textit{categoryScan}, \textit{aiDecision}, dan \textit{manualReview}. Dengan
struktur ini, seluruh hasil inspeksi terhadap satu kontainer dapat ditelusuri
dalam satu entitas logis yang konsisten.

\textbf{Entitas \textit{Manifest}}

Entitas \textit{Manifest} menyimpan data manifes secara lebih lengkap, meliputi
informasi kapal, asal dan tujuan, pihak pengirim dan penerima, deskripsi muatan,
berat, volume, data kepabeanan, serta dokumen pendukung. Keberadaan entitas ini
memungkinkan sistem membandingkan hasil deteksi muatan aktual dengan deklarasi
muatan yang tercatat. Pada implementasinya, data manifes diisikan melalui
aplikasi web, baik melalui entri manual maupun impor data yang tersedia pada
layanan aplikasi.

\textbf{Konteks autentikasi dan akses pengguna}

Selain dua entitas utama tersebut, sistem juga menggunakan data pengguna untuk
mendukung autentikasi, otorisasi, dan peninjauan hasil inspeksi melalui aplikasi
web. Namun, data pengguna tidak menjadi pusat agregasi operasi inspeksi, melainkan
berfungsi sebagai konteks akses terhadap data kontainer dan manifes. Pada sisi
arsitektur, hal ini menunjukkan adanya pemisahan yang jelas antara identitas
pengguna aplikasi dan komunikasi mesin-ke-mesin dari komponen \textit{edge},
tanpa menjadikan mekanisme autentikasi sebagai fokus utama pembahasan.

Di luar jalur operasi utama, realisasi sistem juga sempat memiliki kanal tambahan
untuk data pelacakan dan status aliran. Akan tetapi, kanal tersebut tidak
menjadi bagian inti dari operasi inspeksi kontainer yang dibahas pada tugas akhir
ini, sehingga tidak diposisikan sebagai model data utama pada naskah ini.

\subsection{Pemisahan Penyimpanan Terstruktur dan Artefak Visual}
\label{subsec:pemisahan-penyimpanan}

Realisasi sistem juga menunjukkan bahwa penyimpanan data dilakukan dengan
memisahkan data terstruktur dari artefak visual. Data yang bersifat operasional
dan transaksional disimpan di MongoDB, sedangkan citra hasil OCR, visualisasi
hasil deteksi, dan bukti pendukung lainnya disimpan di Cloudflare R2.

Keputusan ini relevan dengan kebutuhan arsitektural yang telah dirumuskan pada
Bab IV. Data terstruktur membutuhkan dukungan kueri, pembaruan status, dan
pengelolaan konsistensi. Sebaliknya, artefak visual membutuhkan media
penyimpanan objek yang lebih sesuai untuk berkas berukuran besar. Pemisahan ini
meningkatkan efisiensi pengelolaan data sekaligus menjaga struktur informasi
tetap ringkas pada basis data utama.

\subsection{Dukungan terhadap Ketertelusuran dan Audit}
\label{subsec:dukungan-ketertelusuran}

Salah satu tujuan penting dari rancangan arsitektur adalah menyediakan
ketertelusuran proses inspeksi. Pada implementasinya, tujuan tersebut
direalisasikan melalui penyimpanan stempel waktu, identitas perangkat,
identitas kamera, hasil deteksi, ringkasan keputusan, serta artefak visual pada
setiap tahapan penting.

Pendekatan ini memungkinkan sistem menjawab pertanyaan operasional seperti
kontainer apa yang diperiksa, kapan pemeriksaan dilakukan, kamera mana yang
digunakan, hasil apa yang ditemukan, dan bukti visual apa yang mendukung
temuan tersebut. Dengan demikian, struktur implementasi yang terbentuk tidak
hanya mendukung kebutuhan operasional harian, tetapi juga mendukung audit,
verifikasi ulang, dan evaluasi proses secara historis.

\section{Implikasi Implementasi terhadap Kontribusi Arsitektural}
\label{sec:implikasi-implementasi}

Artefak realisasi sistem menunjukkan bahwa rancangan arsitektur yang diusulkan dapat
diterjemahkan ke dalam sistem yang berjalan dengan pemisahan tanggung jawab yang
jelas. Komponen \textit{edge} menangani pengumpulan data dan inferensi awal,
layanan API menangani integrasi dan persistensi, aplikasi web menangani
interaksi pengguna, dan lapisan data menangani penyimpanan serta penyediaan
informasi. Pola realisasi ini memperkuat validitas keputusan arsitektural
yang menekankan modularitas dan integrasi berlapis.

Dari sudut pandang kontribusi tugas akhir, artefak realisasi ini menunjukkan bahwa
perancangan arsitektur tidak berhenti pada level konseptual. Struktur data,
jalur integrasi, dan pembagian komponen yang dirancang menjadi dasar
bagi sistem yang direalisasikan pada proyek pengembangan sistem. Hal tersebut penting karena
menunjukkan bahwa kontribusi tugas akhir ini terletak pada kemampuan merancang
arsitektur yang dapat diimplementasikan, terintegrasi, dan mendukung kebutuhan
ketertelusuran proses inspeksi di lingkungan pelabuhan. Dengan kata lain, nilai
utama bab ini bukan pada rincian kode, melainkan pada keterhubungan antara
rancangan arsitektural dan perilaku sistem yang berhasil direalisasikan.
